\addcontentsline{toc}{chapter}{ABSTRACT}
\chapter*{ABSTRACT}
\thispagestyle{empty}

% \textit{You probably have to follow the structure of this first paragraph.}

The thesis titled \textbf{Student Forum Platform: Enhancing Academic Feedback and Communication in Higher Education} was developed by student \textit{LastName FirstName} as a bachelor's project at the \textit{Technical University of Moldova}. It is written in English and contains an introduction, 4 chapters (Domain Analysis, Requirements Specification, System Design, System Implementation), a list of figures, tables, conclusion, bibliography, and appendices.



In this study, we address the \textbf{problem of limited and inefficient student feedback systems} in higher education, where communication between students and institutions remains infrequent, fragmented, and often anonymous only in form. The research identifies the shortcomings of existing solutions such as Moldova’s SIMU platform, which collects evaluations only once per semester and provides no open, anonymous peer discussion channel.

To solve this, the proposed \textbf{Student Forum Platform} introduces a secure, centralized, and anonymous environment enabling students to exchange opinions, post course feedback, and participate in polls. The system is built using \textbf{Next.js}, \textbf{Supabase}, and \textbf{Drizzle ORM}, ensuring scalability, performance, and type-safe database operations. Authentication and data protection are achieved through \textbf{OAuth 2.0}, \textbf{Multi-Factor Authentication (MFA)}, and \textbf{JWT-based session handling}, complying with modern web security standards.

Initial testing and comparative analysis show improved engagement and transparency compared to traditional feedback systems. By integrating anonymity, verified access, and community-driven features, the platform fosters a culture of open communication in academic environments. The project’s anticipated impact is a \textbf{more responsive, student-centered education ecosystem}, where real-time feedback contributes to continuous institutional improvement.

\textbf{Keywords:} Student forum, feedback platform, academic communication, Supabase, Next.js, Drizzle ORM, OAuth2, MFA, higher education